\section*{Bab \ref{ch:g3:shapes} --- Bangun Datar dan Sudut Siku-siku}

\begin{solution}{exo:g3:shapes:1}
\omterm{def:g3:shapes:rightangle}{Sudut siku-siku} yang khas: pojok papan tulis, pojok selembar kertas,
pojok pintu. Sudut yang bukan siku-siku: bukaan sepasang gunting,
atau jarum jam pada pukul 1. Masing-masing diperiksa dengan
penggaris siku (\cref{met:g3:shapes:check}) --- jangan pernah hanya dengan mata.
\end{solution}

\begin{solution}{exo:g3:shapes:2}
Lukisannya: satu sisi $6$ cm, \omterm{def:g3:shapes:rightangle}{sudut siku-siku} di kedua ujungnya
dengan penggaris siku, sisi $4$ cm, lalu tutuplah bangunnya. Tanda:
$4$ persegi kecil di pojoknya, satu garis kecil pada tiap sisi $6$
cm, dan dua garis kecil pada tiap sisi $4$ cm.
\end{solution}

\begin{solution}{exo:g3:shapes:3}
Persegi punya $4$ \omterm{def:g3:shapes:rightangle}{sudut siku-siku} dan $4$ sisi sama panjang --- semuanya ditandai.
\end{solution}

\begin{solution}{exo:g3:shapes:4}
Lukisan di kertas berpetak; petaknya menjamin \omterm{def:g3:shapes:rightangle}{sudut siku-siku}, dan
membilang petak menjamin panjangnya.
\end{solution}

\begin{solution}{exo:g3:shapes:5}
``Setiap persegi adalah persegi panjang'': \emph{benar} --- persegi
punya $4$ \omterm{def:g3:shapes:rightangle}{sudut siku-siku}, dan hanya itulah syarat persegi panjang.

``Setiap persegi panjang adalah persegi'': \emph{salah} --- persegi
panjang $6 \times 4$ punya sisi yang tidak sama panjang.

``Sebuah segitiga dapat punya dua \omterm{def:g3:shapes:rightangle}{sudut siku-siku}'': \emph{salah}
--- kedua sisi yang ditarik dari dua \omterm{def:g3:shapes:rightangle}{sudut siku-siku} itu akan tegak
lurus pada garis alas yang sama, sehingga tidak pernah bertemu untuk menutup segitiganya.
\end{solution}

\begin{solution}{exo:g3:shapes:6}
$OA = 4$ cm (\omterm{def:g3:shapes:shapes}{jari-jarinya}). $OB = 4$ cm juga: \emph{setiap} titik
pada lingkaran berjarak sama dari pusatnya.
\end{solution}

\begin{solution}{exo:g3:shapes:7}
Persegi panjang itu terpotong menjadi tiga segitiga: dua segitiga
siku-siku ($AMD$ dan $MBC$, siku-siku di $A$ dan di $B$) serta
segitiga $MCD$ di tengah.
\end{solution}

\begin{solution}{exo:g3:shapes:8}
Yang punya sedikitnya satu \omterm{def:g3:shapes:rightangle}{sudut siku-siku}: persegi, persegi
panjang, segitiga siku-siku, dan huruf L (pojok dalam maupun pojok
luarnya). Lingkaran sama sekali tidak punya pojok.
\end{solution}

\begin{solution}{exo:g3:shapes:9}
Setiap sisi menempuh ``$2$ ke kanan, $1$ ke atas'' (atau versi yang
diputar: ``$1$ ke kanan, $2$ ke bawah'', dan seterusnya): setiap
sisi adalah diagonal blok petak $2 \times 1$, jadi keempatnya sama panjang.
\end{solution}

\begin{solution}{exo:g3:shapes:10}
Satu program yang benar: ``1. Gambarlah persegi panjang $ABCD$
dengan $AB = 3$ cm (alas) dan $BC = 2$ cm. 2. Di atas sisi $[DC]$,
gambarlah persegi $DCEF$ dengan sisi $3$ cm, di luar persegi panjang
itu.'' Setiap program yang jelas, bernomor dan lengkap itu benar ---
ujilah dengan mengikutinya kata demi kata.
\end{solution}

\begin{solution}{exo:g3:shapes:11}
Persegi $1 \times 1$: $9$. Persegi $2 \times 2$: $4$.
$3 \times 3$: $1$. Seluruhnya: $9 + 4 + 1 = 14$ persegi.

\end{solution}
